\section{Related Work}
\label{sec:relatedwork}

\subsection{Reproduction and replication studies in software engineering}

Reproduction studies in software engineering typically re-execute a published technique and
compare outcomes, reporting how many results hold. This work does not re-execute anything.
It asks a prior question -- whether the released evidence would let a re-execution be checked
against the original at all -- and answers it per evidence item rather than per study. The
distinction matters because a study can fail a re-execution for reasons that are invisible
when only the aggregate is compared, and can pass one while the correspondence between the
release and the published table remains unestablished.

\subsection{Artifact evaluation and availability badging}

Artifact evaluation at software-engineering venues certifies that an artifact is available,
functional, or reusable. Availability turns on retrievability -- the link resolves and the
contents are what the paper describes. Twenty-five papers in this cohort expose at least one
artifact locator. Twenty-four have a reachable locator disposition, while one paper's
own-artifact deposit was unreachable. We did not apply any venue's badging
rubric, so we do not claim these artifacts would earn a badge; we observe only that the
property those rubrics centre on is one this cohort largely satisfies.

The property it does not satisfy is a complete outcome-specific evidence chain. Thirteen of
the 20 own-artifact papers in this cohort release per-run outputs, while every decision in the
four deep-audit studies lacks at least one outcome-specific public candidate and therefore
cannot establish a complete binding. A
criterion asking for that mapping -- which release, which run set, which evaluator
configuration, which task universe, which aggregation rule -- would separate the two, and
retrievability does not.

\subsection{Benchmark quality and contamination in coding-agent evaluation}

A growing line of work audits the benchmarks these agents are evaluated on: insufficient
test cases that let wrong patches pass, solution leakage in issue text, and contamination of
task sets by model training data. Several such studies appear in our own cohort, since they
are themselves published comparisons of coding agents~\cite{UtboostRigorousEvaluation2025,RevisitingSweBench2025}.

That work asks whether the measuring instrument is sound. This paper asks whether the
measurement's evidence trail is recoverable, and the two are independent: a benchmark can be
flawless and a study's release still fail to say which run produced its table, or a benchmark
can be contaminated in ways only recoverable per-run evidence would expose. Our bounded
released-record checks are of the second kind: AutoCodeRover reports three repetitions of its
base configuration, while its public new- and old-evaluation run reports expose different
unbound candidate counts. This is an
artifact observation, not an end-to-end recomputation.

\subsection{Preregistration and specification drift}

This audit uses dated design records: required-item vectors, the frozen roster, and the planned
two-tier split and sampling seed were recorded before the outcomes they bear on were inspected.
The latter produced a seeded draw but no tier-two compatibility results. We also
report where our own specification moved -- the eligible-decision denominator changed three
times as the eligibility rule was corrected -- because an audit that conceals its own drift
is in no position to report anyone else's.
